\section{Probability Generating Functions}

Note that all natural numbers \(\NN\) in this section includes \(0\).

Let \(X\) be a random variable with values in \(\NN\). We denote the probability mass function of \(X\) as
\[
    p_r = \Prob\pqty{X = r}, \quad r \in \NN.
\]

\subsection{Definitions and Properties}

\begin{definition}[Probability Generating Function]
    \label{def:probability-generating-function}%
    The \emph{probability generating function} of \(X\) is defined to be
    \[
        p \pqty{z} = \Expt\bqty{z^X} = \sum_{r = 0}^{\infty} p_r \cdot z^r = p_0 + p_1 \cdot z + p_2 \cdot z^2 + \cdots
    \]
    where \(\abs{z} < 1\).
\end{definition}

Note that this infinite series is indeed well-defined. Taking the absolute value, we have
\begin{align*}
    \abs{p\pqty{z}} &= \abs{\sum_{r = 0}^{\infty} p_r \cdot z^r} \\
    &\leq \sum_{r = 0}^{\infty} p_r \cdot \abs{z^r} \\
    &\leq \sum_{r = 0}^{\infty} p_r \\
    &= 1,
\end{align*}
so the series converges absolutely for \(\abs{z} \leq 1\), and hence in particular, the radius of the convergence is at least \(1\).

\begin{theorem}[Uniqueness of Probability Generating Function]
    \label{thm:pgf-uniqueness}%
    The probability generating function uniquely determines the distribution of the random variable \(X\).
\end{theorem}

\begin{proof}
    Let \(\pqty{p_r}\) and \(\pqty{q_r}\) be two probability distributions with the same probability generating functions, i.e., for all \(\abs{z} \leq 1\), we have
    \[
        \sum_{r = 0}^{\infty} p_r z^r = \sum_{r = 0}^{\infty} q_r z^r.
    \]

    We want to show \(p_r = q_r\) for all \(r\). If we put in \(z = 0\), then we have \(p_0 = q_0\).

    We continue by strong induction. Suppose that \(p_r = q_r\) for all \(r \leq n\). We want to show that \(p_{n + 1} = q_{n + 1}\). By the induction hypothesis, we have
    \[
        \sum_{r = n + 1}^{\infty} p_r z^r = \sum_{r = n + 1}^{\infty} q_r z^r.
    \]

    Dividing by \(z^{n + 1}\), and set \(z = 0\), we have
    \[
        p_{n + 1} = q_{n + 1}.
    \]
\end{proof}

\begin{proposition}[Probability Generating Function of Sum]
    \label{prop:pgf-sum}%
    Let \(X_1, \ldots, X_n\) be independent random variables, with probability generating functions
    \[
        q_i\pqty{z} = \Expt\bqty{z^{X_i}}.
    \]

    Then, we have
    \[
        p\pqty{z} = q_1\pqty{z} \cdots q_n\pqty{z}.
    \]
\end{proposition}

\begin{proof}
    We have
    \begin{align*}
        p\pqty{z} &= \Expt\bqty{z^{X_1 + \cdots + X_n}} \\
        &= \Expt\bqty{z^{X_1}} \cdots \Expt\bqty{z^{x_n}}
    \end{align*}
    by independence, as desired.
\end{proof}

\subsection{Probability Generating Functions of Standard Distributions}

\begin{example}[Probability Generating Function of Binomial Distribution]
    \label{ex:pgf-binomial}%
    Suppose \(X \sim \Binomial\pqty{n, p}\). We have
    \begin{align*}
        p\pqty{z} = \Expt\bqty{z^X} &= \sum_{k = 0}^{n} z^k \cdot \binom{n}{k} \cdot p^k \cdot \pqty{1 - p}^{n - k} \\
        &= \pqty{pz + 1 - p}^n.
    \end{align*}
\end{example}

\begin{example}[Sum of Independent Binomial Distributions]
    \label{ex:sum-binomial}%
    Suppose \(X \sim \Binomial\pqty{n, p}\) and \(Y \sim \Binomial\pqty{m, p}\), with \(X \indp Y\). We want to find the distribution of \(\pqty{X + Y}\). We have
    \begin{align*}
        \Expt\bqty{z^{X + Y}} &= \Expt\bqty{z^X} \cdot \Expt\bqty{z^Y} \\
        &= \pqty{pz + 1 - p}^{n} \cdot \pqty{pz + 1 - p}^{m} \\
        &= \pqty{pz + 1 - p}^{n + m}.
    \end{align*}

    Therefore, \(X + Y \sim \Binomial\pqty{n + m, p}\).
\end{example}

\begin{example}[Probability Generating Function of Geometric Distribution]
    \label{ex:pgf-geometric}%
    Suppose \(X \sim \Geometric\pqty{p}\), where we use the convention of failures, starting from zero. Then
    \[
        \Expt\bqty{z^X} = \sum_{r = 0}^{\infty} z^r \pqty{1 - p}^r p = \frac{p}{1 - z\pqty{1 - p}}.
    \]
\end{example}

\begin{example}[Probability Generating Function of Poisson Distribution]
    \label{ex:pgf-poisson}%
    Suppose \(X \sim \Poisson\pqty{\lambda}\). Then
    \[
        \Expt\bqty{z^X} = \sum_{r = 0}^{\infty} z^r \cdot e^{-\lambda} \cdot \frac{\lambda^r}{r!} = e^{-\lambda} \cdot e^{\lambda z} = e^{\lambda \pqty{z - 1}}.
    \]
\end{example}

\begin{example}[Sum of Independent Poisson Distributions]
    Suppose \(X \sim \Poisson\pqty{\lambda}\), \(Y \sim \Poisson\pqty{\mu}\), and \(X \indp Y\). Then we have
    \[
        \Expt\bqty{z^{X + Y}} = e^{\lambda \pqty{z - 1}} \cdot e^{\mu\pqty{z - 1}} = e^{\pqty{\lambda + \mu} \pqty{z - 1}},
    \]
    so \(X + Y \sim \Poisson\pqty{\lambda + \mu}\).
\end{example}

\subsection{Expectation and Moments}

We now investigate how to find the expectation (and moments) using the probability generating function.

\begin{theorem}[Expectation from Probability Generating Function]
    \label{thm:expt-pgf}%
    We have
    \[
        \lim_{z \to 1^{-}} p'\pqty{z} = p'\pqty{1^{-}} = \Expt\bqty{X}.
    \]
\end{theorem}

\begin{proof}
    Suppose \(\Expt\bqty{X} < \infty\). Take \(0 < z < 1\), so
    \[
        p'\pqty{z} = \sum_{r = 1}^{\infty} r p_r z^{r - 1} \leq \sum_{r = 1}^{\infty} r p_r = \Expt\bqty{X}.
    \]

    Therefore, \(p'\pqty{z}\) is an increasing function of \(z\), and hence,
    \[
        \lim_{z \to 1^{-}} p'\pqty{z} \leq \Expt\bqty{X}.
    \]

    Let \(\epsilon > 0\). There exists \(N\) large, such that
    \[
        \sum_{r = 1}^{N} r p_r \geq \Expt\bqty{X} - \epsilon.
    \]

    Therefore,
    \[
        \lim_{z \to 1^{-}} p'\pqty{z} \geq \lim_{z \to 1^{-}} \sum_{r = 1}^{N} r p_r z^{r - 1} = \sum_{r = 1}^{N} r p_r \geq \Expt\bqty{X} - \epsilon.
    \]

    This finishes the proof.

    Now, suppose \(\Expt\bqty{X} = \infty\). Then for any \(M > 0\), there is some \(N \in \NN\), such that
    \[
        \sum_{r = 1}^{N} r p_r \geq M.
    \]

    Therefore, from before, we have
    \[
        \lim_{z \to 1^{-}} p'\pqty{z} \geq \lim_{z \to 1^{-}} \sum_{r = 1}^{N} r p_r z^{r - 1} = \sum_{r = 1}^{N} r p_r \geq M.
    \]

    Therefore, \(\lim_{z \to 1^{-}} p'\pqty{z} = \infty\).
\end{proof}

Similarly, we may find the factorial moments by taking the derivatives.

\begin{theorem}[Factorial Moments from Probability Generating Function]
    \label{thm:factorial-moments-pgf}%
    For the second factorial moment, we have
    \[
        p''\pqty{1^-} = \lim_{z \to 1^-} p''\pqty{z} = \Expt\bqty{X \pqty{X - 1}}.
    \]

    In general, for any \(k \geq 0\), we have
    \[
        p^{\pqty{k}} \pqty{1^-} = \lim_{z \to 1^-} p^{\pqty{k}}\pqty{z} = \Expt\bqty{X \pqty{X - 1} \cdots \pqty{X - k + 1}}.
    \]
\end{theorem}

\begin{corollary}[Variance from Probability Generating Function]
    \label{cor:var-pgf}%
    We have
    \[
        \Var\pqty{X} = p''\pqty{1^-} + p'\pqty{1^-} - \pqty{p'\pqty{1^-}}^2.
    \]
\end{corollary}

\begin{proposition}[Probability from Probability Generating Function]
    \label{prop:prob-pgf}%
    We have
    \[
        \Prob\pqty{X = n} = p^{\pqty{n}} \pqty{0} \cdot \frac{1}{n!}.
    \]
\end{proposition}

\subsection{Sum of a Random Number of Random Variables}

Let \(X_1, X_2, \ldots\) be independent and identically distributed, and let \(N\) be a random variable (independent) and take values in \(\NN\).

We set the sum
\[
    S_n = X_1 + \cdots + X_n = \sum_{i = 1}^{n} X_i,
\]
and as a function,
\[
    S_n\pqty{\omega} = \sum_{i = 1}^{n} X_i\pqty{\omega}.
\]

\begin{definition}[Sum of a Random Number of Random Variables]
    \label{def:sum-random-number-rvs}%
    We define
    \[
        S_N \pqty{\omega} = S_{N\pqty{\omega}} \pqty{\omega} = \sum_{i = 1}^{N\pqty{\omega}} X_i\pqty{\omega}.
    \]
\end{definition}

\begin{lemma}[Probability Generating Function of a Sum]
    \label{lem:pgf-sum}%
    Let \(q\pqty{z} = \Expt\bqty{z^N}\) be the probability generating function of \(N\), and \(p\pqty{z} = \Expt\bqty{z^{X_i}}\) be the probability generating function of \(X_i\)s. Then we have
    \[
        r\pqty{z} = \Expt\bqty{z^{S_N}} = q\pqty{p\pqty{z}}.
    \]
\end{lemma}

\begin{proof}[\proofname\ by Indicator Function]
    We let \(r\pqty{z} = \Expt\bqty{z^{S_N}}\). We have
    \begin{align*}
        r\pqty{z} &= \Expt\bqty{z^{S_N}} \\
        &= \Expt\bqty{z^{S_N} \cdot \sum_{n = 0}^{\infty}  \indicator\pqty{N = n}} \\
        &= \sum_{n = 0}^{\infty} \Expt\bqty{z^{S_n} \cdot \indicator \pqty{N = n}} \\
        &= \sum_{n = 0}^{\infty} \Expt\bqty{z^{X_1 + \cdots + X_n} \cdot \indicator\pqty{N = n}} \\
        &= \sum_{n = 0}^{\infty} \Expt\bqty{z^{X_1 + \cdots + X_n}} \cdot \Prob\pqty{N = n} \\
        &= \sum_{n = 0}^{\infty} \pqty{p\pqty{z}}^{n} \cdot \Prob\pqty{N = n} \\
        &= q\pqty{p\pqty{z}}.
    \end{align*}
\end{proof}

\begin{proof}[\proofname\ by Conditional Expectation]
    We have
    \[
        r\pqty{z} = \Expt\bqty{z^{S_N}} = \Expt\bqty{\Expt\bqty{\Conditional{z^{S_N}}{N}}} = \Expt\bqty{g\pqty{N}},
    \]
    where
    \begin{align*}
        g\pqty{n} &= \Expt\bqty{\Conditional{z^{S_N}}{N=n}} \\
        &= \Expt\bqty{\Conditional{z^{S_n}}{N = n}} \\
        &= \Expt\bqty{z^{S_n}} \\
        &= \pqty{p\pqty{z}}^{n}.
    \end{align*}

    Hence, we have
    \[
        r\pqty{z} = \Expt\bqty{p\pqty{z}^{N}} = q\pqty{p\pqty{z}}.
    \]
\end{proof}

\begin{corollary}[Expectation of a Random Sum]
    \label{cor:expt-random-sum}%
    We have
    \[
        \Expt\bqty{S_N} = \Expt\bqty{N} \cdot \Expt\bqty{X_i}.
    \]
\end{corollary}

\begin{proof}
    We have
    \[
        \Expt\bqty{S_N} = r' \pqty{1^-} = q' \pqty{p\pqty{1^-}} \cdot p'\pqty{1^-} = \Expt\bqty{N} \cdot \Expt\bqty{X_i}.
    \]
\end{proof}

\begin{corollary}[Variance of a Random Sum]
    \label{cor:var-random-sum}%
    We have
    \[
        \Var\pqty{S_N} = \Expt\bqty{N} \cdot \Var\pqty{X_i} + \Var\pqty{N} \cdot \pqty{\Expt\bqty{X_i}}^2.
    \]
\end{corollary}

\begin{proof}
    By differentiating the probability generating function twice.
\end{proof}